\section{Derivations}
\label{app:derivations}

\subsection{Reproduction number and suppression efficacy}
\note{Next-generation matrix of the $(I, V)$ subsystem at the disease-free state; $R_0 = \frac{p(1-\mu)}{\delta}\,\frac{\beta T^*}{c + \beta T^*}$; solving $R_0(\varepsilon) = 1$ for an infection-blocking drug gives $1 - \varepsilon^* = \delta c / (\beta T^* (p(1-\mu) - \delta))$.}

\subsection{The activation threshold as the chronic set point}
\note{Any equilibrium with $E > 0$ has $q(I_w + I_m) = d_E$.}

\subsection{Establishment with fixed effectors}
\note{Generating functions of the two-type process; the fixed-point equations; the one-type closed form $\pi(1 - 1/R_0)$; vectorised closed form for the mutant.}

\subsection{Lineages in a changing environment}
\note{Kendall's formula $\rho(t) = \int_0^t (d - b)$, $W(t) = 1 + \int_0^t d\, e^{\rho}$, $p_0(t) = 1 - 1/W(t)$; the effector lineage with $b = q I(t)$ and $d = d_E$; the mutant lineage with lumped rates $b = p_m \pi_m(T(t))$, $d = \delta + k E(t)$, and the check that the lumped process reproduces $1/R$ in a constant environment; the Poisson de novo term and conditioning on the wild type's survival time; the decaying-effector version used for the regime map.}
